\chapter{ICT Law}

\section{ICT Act 2006}

\subsection{Section-68: Establishment of Cyber Tribunal}
\begin{itemize}[parsep=0.5em]
	\item	The Government shall, by notification in the Official Gazette, establish one or more Cyber Tribunals for speedy and effective trials of offences under this Act.
	\item	The Tribunal,
	established under sub-section (1),
   % in consultation with the Supreme Court,
   shall be constituted by a Session Judge or Additional Session Judge appointed by the Government as “Judge, Cyber Tribunal.”
	\item	The Tribunal may have jurisdiction over entire Bangladesh or one/more Session Divisions and prosecutes only offences under this Act.
	\item	Ongoing cases in any Session Court are not automatically transferred to the Tribunal; the Government may transfer cases by notification in the Official Gazette.
	\item	The Tribunal need not retake witness statements, rehear, or restart activities already undertaken, but shall continue prosecution from existing records.
	\item	The Government shall define the place and time for the Tribunal
   s activities, which it shall conduct accordingly.
\end{itemize}


\subsection{Section-69: Trial Procedure of Cyber Tribunal}
\begin{itemize}[parsep=0.5em]
	\item	The Tribunal shall not accept any offence trial without a written report from a police officer not below Sub-Inspector or prior approval of the Controller or an authorized officer.
	\item	The Tribunal shall follow the rules mentioned in Chapter 23 of the Code of Criminal Procedure, if not inconsistent with this Act.
	\item	Prosecution cannot be suspended without written reasons and for just adjudication.
	\item	If the accused has absconded and cannot be arrested immediately, the Tribunal may order appearance via publication in two national Bengali dailies; failure to appear allows prosecution in absence.
	\item	Sub-section (4) does not apply if the accused absconds after getting bail.
	\item	The Tribunal can order police, the Controller, or an authorized officer to reinvestigate and submit a report within a stipulated time, on its own initiative or on application.
\end{itemize}

\pagebreak
\subsection{Section-71: Rules Relating to Bail}
The Judge of Cyber Tribunal shall not grant bail to any person accused under this Act unless:  

\begin{enumerate}[itemsep=0.5em]
    \item The Government is given a hearing opportunity on similar bail orders.
    \item The Judge is satisfied that:
    \begin{enumerate}
        \item There is reasonable cause to believe the accused may not be proved guilty.
        \item The offence is not severe and the punishment shall not tough even if guilt is proved.
    \end{enumerate}
    \item The Judge records the reasons for such satisfaction.
\end{enumerate}

\vspace{2em}
\section{Importance of ICT Law of CSE}

ICT law is crucial for CSE students as it provides the legal framework for digital work, ensuring ethical and responsible technology development and usage.
% It helps students handle legal challenges in software, data, and online interactions, fostering responsible innovation and preventing legal issues.
Here's why it's important:

% Importance of ICT Law for CSE Students:

\begin{enumerate}[itemsep=1em, leftmargin=1em]
   \item \textbf{Navigating the Digital Legal Landscape}
   \begin{itemize}[itemsep=0.5em]
      \item	\textbf{Software and Intellectual Property:} Students must understand copyright, patents, and trade secrets to protect their creations and respect others' rights.
      \item	\textbf{Data Protection and Privacy:} Students need to be
      aware of laws like GDPR and CCPA for user privacy and data security.
      \item	\textbf{Cybersecurity and Cybercrimes:} Understanding cyber laws helps develop secure systems and prevent attacks.
      \item	\textbf{E-commerce and Transactions:} Students need awareness of legal aspects of e-commerce, contracts, and digital signatures.
      \item	\textbf{Ethical Hacking and Penetration Testing:} ICT law ensures these practices are conducted within legal boundaries.
   \end{itemize}
   
   \item \textbf{Promoting Responsible Innovation}
   \begin{itemize}[itemsep=0.5em]
      \item	\textbf{Developing Secure and Ethical Technology:} Legal awareness ensures technologies are functional, ethical, and compliant.
      \item	\textbf{Preventing Legal Issues:} Helps avoid pitfalls such as copyright infringement, data breaches, and unfair competition.
      \item	\textbf{Fostering Innovation:} Knowledge of legal frameworks empowers responsible innovation within ethical and legal limits.
   \end{itemize}
   
   \item \textbf{Career Opportunities}
   \begin{itemize}[itemsep=0.5em]
      \item	\textbf{Specialized Roles:} Growing demand in roles like cybersecurity lawyers, data protection officers, and legal tech consultants.
      \item	\textbf{Interdisciplinary Careers:} ICT law knowledge benefits careers combining computer science with law, policy, or business.
   \end{itemize}
\end{enumerate}

In essence ICT law provides the foundation for responsible, secure, and ethical technology development, enabling CSE students to contribute to a safer digital world.
